\section{Evaluation}
\label{sec:evaluation}

We evaluate \sys and baselines on three \ac{MoE} models and ask the following questions.
First, how much fine-tuning throughput does \sys obtain and what is the impact on inference latency, compared with an iteration-level colocation system (\llmstation) and with uncoordinated colocation (\S\ref{sec:eval-e2e})?
Second, where does the throughput in \sys come from, and what execution policies are needed to leverage
intra-layer idle time safely (\S\ref{sec:eval-ablation}, \S\ref{sec:eval-latency})?
Third, does \sys still provide  benefits once the serving engine balances expert load itself (\S\ref{sec:eval-eplb})?
Appendix~\ref{sec:appx-eval} details our experimental setup, and Appendix~\ref{sec:appx-experiments} reports more extensive experiments: all models and prompt lengths, fewer GPUs, a second production trace, the \uncoord \ac{SM}-cap sweep, the choice of design parameters, bubble utilization, and overheads. Appendix~\ref{sec:appx-gap} isolates the value of per-GPU planning with a trace-driven replay.
%\S\ref{sec:appx-gap} separately isolates two prerequisites for bubble filling: router-aware per-GPU planning and sufficiently fine-grained training work. 
%The experiments below evaluate their end-to-end use by the complete system.

\subsection{Experimental setup}
\label{sec:eval-method}

% Here, we provide the most important details of our experimental setup.
% \amk{that sentence above can go if we need space}

\paragraph{Hardware and models.}
We run all experiments on a single compute node equipped with four \textsc{NVLink}-connected A100-SXM4-80GB GPUs, serving with \vllm in BF16 precision.
We evaluate three \ac{MoE} models: \textsc{Qwen3-30B-A3B}~\cite{qwen3}, which has only routed experts, 
\textsc{Qwen1.5-MoE-A2.7B}~\cite{qwen1.5_moe} and \textsc{DeepSeek-V2-Lite}~\cite{deepseekv2}, which additionally have shared experts that process every token and therefore add the same work to every GPU.
%Every system fine-tunes the same LoRA job on alpaca-cleaned \oana{Alpaca?} samples.
%\sys uses an early-close margin of $m=\SI{60}{\micro\second}$ (\S\ref{sec:design-optC}) and opens bubble windows only in prefill steps of at least 4k tokens summed across prompts in a batch (\S\ref{sec:design-predict}).

\paragraph{Workloads and metrics.}
For inference, we use \textsc{arXiv} summarization prompts~\cite{cohan-etal-2018-discourse} with lengths of 1k--32k tokens and generate 32 output tokens for each prompt.
We set \vllm's maximum number of sequences to 256.
For this controlled workload, we determine three load regimes.
With a \emph{high} load, the arrival rate is just below the point at which serving alone starts queueing.
Medium and low load use lower rates, which we specify per experiment and workload in the appropriate plots.
We also replay two production traces, \emph{code} and \emph{conversation} from the Azure LLM inference trace~\cite{splitwise}, each over a 10-minute window with the trace's own arrival times.
For fine-tuning, we use the \alpaca dataset, consisting of \num{52000} instruction-demonstration prompts~\cite{alpaca}. %, truncated or padded to \num{128} tokens.

% \paragraph{Baselines.}
% We use LLMStation as the primary colocation baseline because it targets
% the headroom available in our regime. At our prompt lengths, prefill
% already saturates the GPU~\cite{sarathi}, leaving little of the prefill
% compute headroom targeted by systems such as
% DeltaServe~\cite{deltaserve} and FlexLLM~\cite{flexllm}.
% The remaining serving-level opportunities are therefore primarily
% decode and idle-time gaps, which LLMStation is designed to exploit.

\paragraph{Baselines.}
We compare against \llmstation~\cite{llmstation}, a state-of-the-art system for colocating fine-tuning with \ac{LLM} inference on shared GPUs.
\llmstation transforms a fine-tuning job into a suspendable pipeline and interleaves it with inference at iteration granularity, batching tuning work alongside decode steps.
The system leverages the compute left unused by memory-bound decode iterations, and the idle periods in which no inference work is pending,  which makes it the strongest representative of serving-level colocation.
\flexllm~\cite{flexllm} and \deltaserve~\cite{deltaserve} instead place fine-tuning work alongside prefill, but only when their per-iteration planner sees headroom in the serving iteration. At our prompt lengths, prefill saturates GPU compute~\cite{sarathi}, and an iteration-level view cannot see the per-GPU slack inside \ac{MoE} layers (\S\ref{sec:gap}). Their planners therefore admit almost no work into prefill. Since their training progress is tied to these admissions, they also make little use of decode slack, which leaves them mostly the idle time that \llmstation already exploits.
This makes iteration-level colocation by \llmstation the more appropriate point of comparison.
We also use \emph{\Uncoord} as a baseline, which runs a stock \ac{LoRA} trainer beside \vllm.

\paragraph{Implementation.}
We implement \sys in Python in about 8k lines on \vllm~0.21.1 (v1 engine), \textsc{PyTorch}~2.11, and \textsc{Triton}~3.6.
The offline compiler builds on \texttt{AOTAutograd} and \texttt{Inductor}.
At runtime, hooks in \vllm's \ac{MoE} layer predict and open each window.
A separate trainer process per GPU runs the compiled chunks under CUDA MPS and accesses the serving weights through CUDA IPC, so they are not duplicated.

For \llmstation, we port its artifact to our vLLM/\textsc{PyTorch} versions while preserving its scheduling policy.
Because its released fusion engine assumes dense projections with fixed weight placement, which does not support expert-parallel \ac{MoE}, we use its separate \textsc{PyTorch} training path under CUDA MPS.
We evaluate several \ac{MoE} expert-execution implementations and report the fastest, using padded batched \acp{GEMM} through \textsc{cuBLAS}.
\llmstation's latency is normalized to the same serving engine with training disabled.
\emph{\Uncoord} runs a stock torchtune \ac{LoRA} trainer beside \vllm under
CUDA MPS, with its \ac{SM} share capped at 40\%.

\paragraph{Metrics.}
For fine-tuning, we report throughput in samples/s.
For serving, we report p90 \ac{TTFT} and \ac{TPOT}, normalized to serving alone, which capture prefill and per-token decode latency, respectively.
Each experiment is repeated three times and we average all results.
% \TODO{(error bars show standard deviation)}.

\begin{figure}[t]
  \centering
  \inputplot{ft_throughput_qwen3}
  \caption{Fine-tuning throughput for \sys, \llmstation, and \uncoord (40\% cap) on \textsc{Qwen3-30B-A3B} across prompt lengths and serving loads. The labels above bars report \sys's fine-tuning throughput relative to \llmstation.}
  \label{fig:eval-e2e-tput}
\end{figure}

\subsection{End-to-end performance of \sys and baselines}
\label{sec:eval-e2e}

We measure end-to-end performance of \sys and baselines on the \textsc{Qwen3-30B-A3B} \ac{LLM}.

% We evaluate all three models over prompt lengths from 1k to 32k tokens
% and at low, medium, and high serving load.  //oana I cut this because it's repetitive with the methodology section.
% For space, the main text shows Qwen3-30B-A3B at 4k, 16k, and 32k.
% Appendix~\ref{sec:appx-e2e} reports Qwen1.5-MoE-A2.7B and DeepSeek-V2-Lite, as well as 1k, 2k, 8k, and 24k prompts.
% Their results are very similar to those below, with the exception of \ac{TPOT} on Qwen1.5-MoE at long prompts (discussed in Appendix~\ref{sec:appx-e2e}).

% Uncoordinated 40%-cap Qwen3 32k@0.5 and 2k@12 rerun 2026-10-27
% (BubbleTea-results/naive_torchtune/runs_cap40); the earlier runs had failed.

\paragraph{Steady-state workloads.}
\Cref{fig:eval-e2e-tput} shows fine-tuning throughput of \sys and baselines, for varying serving loads and 4k--32k prompt lengths.
Across all settings, \sys provides more fine-tuning throughput than \llmstation and trains 2.7--5.0$\times$ faster than \llmstation, with a median improvement of 2.9$\times$ (1.9--6.2$\times$ on all models, \S\ref{sec:appx-e2e}).
%\mv{One sentence here on the uncoordinated baseline and some interesting finding.}
%\rnp{}
\Uncoord obtains a throughput of just 0.26--0.43 samples/s, 2.7--12.4$\times$ less than \sys, and lifting its \ac{SM} cap does not help: the off-the-shelf trainer is bound by its own per-expert kernel launches.
%Across all three models and the full 1k--32k evaluation, \sys{} ,
While \sys and \llmstation both fine-tune during decode and idle time, \sys uses fine-grained compiled chunks, which enables more efficient scheduling. 
%\S\ref{sec:eval-ablation} separates this execution benefit from the additional throughput provided by bubble filling.

\begin{figure}[t]
  \centering
  \inputplot{e2e_latency_qwen3}
  \caption{p90 \ac{TTFT} (top) and p90 \ac{TPOT} (bottom) ratio, relative to the corresponding serving-only execution, for \sys, \llmstation, and \uncoord (40\% cap) on \textsc{Qwen3-30B-A3B} across prompt lengths and serving loads. The dashed line marks serving alone and the dotted line a 10\% increase.}
  \label{fig:eval-e2e-lat}
\end{figure}

\Cref{fig:eval-e2e-lat} shows the p90 \ac{TTFT} and p90 \ac{TPOT} in the same settings as in \Cref{fig:eval-e2e-tput}, relative to a serving-only setting.
Across all 63 configurations, \sys keeps both p90 \ac{TTFT} and p90 \ac{TPOT} within 10\% of serving alone in 50 configurations (79\%), compared with 44 (70\%) for \llmstation.
The median increase of \sys is 4\% and 6\% for p90 \ac{TTFT} and p90 \ac{TPOT}, respectively.
%\mv{Again, an interesting sentence here about the uncoordinated baseline + reference to \Cref{sec:appx-naive}}
%\rnp{Draft: }
\emph{\Uncoord} meets this 10\% bound in only 22 of the 63 configurations: it slows every serving step by a similar amount, which raises p90 latency by 10--20\% over serving alone when serving has slack, but at high load pushes serving past its saturation point, so requests queue and p90 \ac{TTFT} reaches 6--12$\times$ that of serving alone.
Thus, we conclude that \sys across most configurations achieves higher fine-tuning throughput without systematically increasing serving latency.
% As a diagnostic, replacing these exchanges with
% no-ops returns decode latency to the serving-only baseline, identifying
% training communication rather than bubble admission as the source of
% the interference. We report the complete per-workload results and this
% diagnostic in the appendix.

% At 1k and 2k tokens, prefills are too short for \sys{} to expose useful
% bubble windows. In this regime, \sys{} effectively reduces to its
% two-chain decode/idle-time colocation path, and its gains therefore
% measure the efficiency of that path rather than bubble filling.
% The appendix reports these short-prompt results together with a
% breakdown of the available idle, decode, and bubble capacity across
% loads.

%\Uncoord generally either trains far less or inflates serving latency; Appendix~\ref{sec:appx-naive} shows how this tradeoff moves with the trainer's SM cap. \oana{cut if need space}



\paragraph{Production traces.}
We additionally replay two production traces with bursty, time-varying arrivals.
\sys shows 3.2--3.3$\times$ higher fine-tuning throughput than \llmstation while keeping p90 \ac{TTFT} and \ac{TPOT} within 7\% of serving alone.
\Cref{fig:intro-trace} in \S\ref{sec:intro} shows one representative trace.
Thus, even on production traces, \sys achieves superior fine-tuning throughput with low effect on inference latency.

% \oana{overall, this subsection reads a lot like a laundry list of results, without understanding why the results are this way. The following sections give more of an intuition about the mechanisms that lead to the results. You need to emphasize this here as well, especially for \sys. }
% \mv{I agree; we have to give more intuitions on why results are like this.}
% \rnp{Draft paragraph, \textbf{Why \sys behaves this way.} \sys's advantage has two sources, which \S\ref{sec:eval-ablation} separates. At low and medium load, most training still runs in decode steps and idle time, where \sys's short compiled chunks use the headroom more efficiently than \llmstation's layer-granularity tasklets. At high load, idle time disappears and decode steps become scarce, so \llmstation's throughput drops to 0.20--1.00 samples/s. Prefill bubbles, in contrast, grow with the prefill work, so \sys still trains at 1.0--2.8 samples/s and its advantage rises to 2.8--5.0$\times$ (\Cref{fig:eval-e2e-tput}). Latency stays close to serving alone because each window is bounded by the router's prediction and closes before the bubble ends, so training rarely runs on a GPU's critical path.}

\begin{figure}[t]
  \centering
  \inputplot{sources_8k}
  % \caption{
  %   \textbf{Sources of \sys{}'s fine-tuning throughput} (8k prompts).
  %   \cfgColo uses conventional decode and idle-time headroom;
  %   \cfgOptColo additionally keeps training work ready across
  %   dependencies; \sys{} additionally fills intra-layer prefill
  %   bubbles. Bubble filling becomes more valuable as serving load
  %   consumes conventional headroom.\mv{Needs rewording}
  % }
  \caption{Fine-tuning throughput for \emph{decode colocation} (coarse-grained work scheduling), \sys-\textsc{Flat} (fine-grained work scheduling but not within bubbles), and \sys on \textsc{Qwen3-30B-A3B} with prompts of 8k length across serving loads.}
  \label{fig:eval-sources}
\end{figure}

\subsection{Where does \sys's throughput come from?}
\label{sec:eval-ablation}

% Decode colocation (DC) = LLMStation-like colocation (leveraging decode + idle time in prefill)
% Optimizized colocation (no bubble scheduling)
% Soda

Next, we investigate where the fine-tuning throughput in \sys comes from.
\Cref{fig:eval-sources} shows fine-tuning throughput for varying serving loads, and for different colocation mechanisms, while using prompts with length 8k.
We first compare \emph{\cfgColo}, which uses \llmstation-like scheduling at the granularity of an iteration, with a version of \sys, named \sys-\textsc{Flat}, that schedules work at finer granularity but does not schedule work in bubbles.
This comparison  isolates the benefit of \sys's fine-grained scheduling while holding the available serving capacity fixed to conventional decode and idle-time headroom.

% \Cref{fig:eval-sources} shows that \emph{decode colocation (DC)} achieves a throughput of 3.2 samples/s with low serving load, which decreases to 0.8 with a high load.

\sys-\textsc{Flat} improves fine-tuning throughput by 36--60\% over \emph{\cfgColo}, with improvements more pronounced in the low load regime.
This highlights that \emph{the fine-grained scheduling mechanism in \sys is able to significantly increase fine-tuning throughput compared to conventional approaches.} \sys-\textsc{Flat} keeps p90 \ac{TTFT} and \ac{TPOT} within 10\% of serving alone in most configurations.

%at low and medium serving load, decode and idle periods still expose
%substantial capacity, but a single fine-tuning chain can leave it unused
%while waiting on cross-GPU dependencies. By keeping work queued on the
%device and running a second independent chain when needed,
%\emph{Thus, the fine-grained scheduling improves throughput compared to conventional colocation.}

% We decompose \sys{}'s throughput gain in two steps.
% First, we compare 
% Second, we compare \cfgOptColo with \sys to isolate the additional capacity provided by prefill bubbles while holding the training runtime fixed.
% In short, \cfgColo$\rightarrow$\cfgOptColo measures the throughput gained by \emph{more efficient use of conventional headroom}, while \cfgOptColo$\rightarrow$\sys{} measures the throughput gained through \emph{bubble filling}.

% Figure~\ref{fig:eval-sources} shows 8k prompts; the ranges below cover
% 4k, 8k, and 16k prompts, with full results in
% Appendix~.

Next, we compare the fine-tuning throughput of \sys-\textsc{Flat} with that of \sys that additionally schedules fine-tuning work in bubbles.
\Cref{fig:eval-sources} shows that at a medium and high load, \sys achieves 24\% and 27\% higher throughput, respectively, compared to \sys-\textsc{Flat}.
This is because with higher load, there are more abundant bubbles during prefill which \sys effectively leverages.
At low load, where conventional headroom is abundant, leveraging bubbles does not help, and slightly diminishes throughput by 3\%.
We attribute this decrease to the additional overhead required by the \sys scheduling mechanism.
%\Cref{fig:eval-sources} shows that as serving load increases, idle gaps shrink and the server spends more time in prefill, leaving progressively less conventional decode and idle-time headroom (Appendix~\ref{sec:appx-capacity}).
%Prefill bubbles, in contrast, occur inside those increasingly frequent prefills and are
%invisible to an iteration-level scheduler. At high load, filling
%them adds a further 27--37\% fine-tuning throughput over
%\cfgOptColo. 
%\emph{Bubbles matter as conventional headroom disappears.}

% Thus, we conclude that \sys obtains higher fine-tuning throughput in medium- and high-loaded regimes.
% In summary, we conclude that the gains in \sys both come from our finer-grained scheduling as well as the opportunity to leverage bubbles during the prefill phase.
%In summary, the complete-system gain therefore reflects both effects: the runtime required for bubble-scale execution uses conventional headroom more efficiently, while router-aware bubble filling adds capacity when that conventional headroom becomes scarce.

\subsection{How does \sys protect serving latency?}
\label{sec:eval-latency}

% \rnp{To make this easier, I'm going to give u a tldr - 
% we take Soda without bubbles (Juice),
% and put things in bubbles naively (that is unthrottled bubbles),
% then, we gate the number of chunks admitted into the bubbles, thats throttled
% notice that nothing is being said about which GPU stuff is being placed on. We finally add inteligence around not placing work on the GPU that's on the critical path. And that's Soda!
% }

The \sys scheduler should not increase inference latency.
To better understand how our scheduler protects this latency, we progressively enable subcomponents of our scheduler and report in \Cref{fig:eval-latency-ablation} the fine-tuning throughput (left), p90 \ac{TTFT} (middle) and p90 \ac{TPOT} (right) for varying serving loads and scheduler configurations, when using prompts with length 4k.
We highlight the results of prompt lengths of 4k because the effects of these different optimizations are most pronounced at this prompt length, making it easiest to distinguish their individual contributions.

\emph{Unconstrained scheduling} admits fine-tuning work on every GPU during idle time until that GPU has completed the all-reduce operation.
This maximizes the amount of training exposed inside prefill, but can compete for compute needed by serving and can run training on the most-loaded GPU that determines the layer's critical path.

\begin{figure}[t]
  \centering
  \inputplot{latency_ablation_4k}
  \caption{Fine-tuning throughput (left), p90 \ac{TTFT} (middle) and p90 \ac{TPOT} (right) relative to serving alone, for \sys-\textsc{Flat}, \emph{unconstrained scheduling}, \emph{constrained scheduling}, and \sys on \textsc{Qwen3-30B-A3B} with 4k prompts across serving loads. The dashed line marks serving alone and the dotted line a 10\% increase.}
  \label{fig:eval-latency-ablation}
\end{figure}

 \emph{Constrained scheduling} adds four mechanisms that protect serving latency. First, waiting for admission becomes a stream-ordered device wait that consumes no GPU compute.
Second, the two fine-tuning chains run concurrently only while prefill takes less than 70\% of recent serving time and are serialized otherwise (\S\ref{sec:design-chunk}).
Third, prefills shorter than 4k tokens are skipped because their bubbles are too small to justify admission.
Fourth, each GPU stops admitting new chunks a margin $m$ before its bubble ends (\emph{early close}), so that no running chunk delays the all-reduce (\S\ref{sec:design-optC}).

\Cref{fig:eval-latency-ablation} shows that at a high load, \emph{unconstrained scheduling} increases fine-tuning throughput from 0.73 to 2.11 samples/s compared to \sys-\textsc{Flat}, but p90 \ac{TTFT} increases to 5.03$\times$ serving alone and p90 \ac{TPOT} to 1.19$\times$.
Fine-tuning that reaches the serving-critical GPU extends the prefill and can queue subsequent requests.
Exposing intra-layer capacity is therefore not sufficient; admission must avoid interfering with the critical path.
\emph{Constrained scheduling} reduces p90 \ac{TTFT} to 1.03$\times$ and p90 \ac{TPOT} to 1.13$\times$, while retaining 1.33 samples/s.
%Compared to \emph{unconstrained scheduling}, \emph{constrained scheduling} reduces p90 \ac{TTFT} from 5.03$\times$ to 1.03$\times$ and p90 \ac{TPOT} from 1.19$\times$ to 1.13$\times$, while retaining 1.33 samples/s of fine-tuning throughput at a high load.
Therefore, \emph{our scheduling constraints are key to protect serving latency.}

Constrained scheduling, however, can still admit fine-tuning work on the most-loaded GPU because it does not distinguish GPUs by their predicted slack.
\sys uses the router-derived budget to prevent that GPU from admitting work and to stop other GPUs before their predicted bubbles end.
\Cref{fig:eval-latency-ablation} shows that this reduces p90 \ac{TTFT} and \ac{TPOT} further to 1.01$\times$ and
1.04$\times$ compared to serving only, respectively, at 1.00 samples/s of fine-tuning throughput.
%The same tradeoff holds across loads and at 8k and 16k prompts (Appendix~\ref{sec:appx-ablation}).

% \oana{remove paragraph below because you show 4k prompts not 8k.}
% These protections deliberately leave bubble capacity unused.
% At 8k prompts and low load, \sys performs fine-tuning during 36\% of the available bubble duration, losing most of the rest to the early-close margin or to fine-tuning work that is not yet ready.
% In return, 98\% of admission windows close before the bubble ends.

%\sys{} adds the final protection. It uses the router-derived per-GPU
%budget to identify which GPUs are predicted to finish early, keeps the
%predicted straggler closed, and stops admission before each predicted
%bubble ends. Figure~\ref{fig:eval-latency-ablation} shows how these
%progressively stronger protections trade fine-tuning throughput for
%serving isolation.

%\paragraph{Aggressive prefill training exposes capacity but harms serving.}

%\paragraph{Removing avoidable contention recovers most latency.}


%\paragraph{Router-bounded windows protect the remaining critical path.}


% These protections intentionally leave some bubble capacity unused.
% At 8k prompts and 2 requests/s, \sys{} trains during 36\% of available
% bubble time, with most of the remainder lost to unavailable training
% work or the early-close margin; 98\% of windows nevertheless close
% before the serving dependency clears
% (Appendix~\ref{sec:appx-bubble-use}). Appendix \S\ref{sec:appx-gap} isolates
% the value of router-aware per-GPU planning, while
% Appendix~\ref{sec:appx-margin} evaluates the early-close margin.

\begin{figure}[t]
  \centering
  \inputplot{eplb_8k}
  \caption{Fine-tuning throughput (left), p90 \ac{TTFT} (middle) and p90 \ac{TPOT} (right) on \textsc{Qwen3-30B-A3B} with prompts of 8k length at medium (2 requests/s) and high (4 requests/s) serving load. Left: each bar shows the throughput of \sys-\textsc{Flat}, with the additional throughput of \sys from bubble filling stacked on top; labels give the gain from bubble filling. Middle and right: latency of serving alone without and with \ac{EPLB}, and of  \sys with \ac{EPLB}.}
  \label{fig:eval-eplb}
\end{figure}

\subsection{Composition with load balancing optimizations}
\label{sec:eval-eplb}

We next ask whether \sys remains useful when serving already uses an expert load-balancing optimization.
We use \vllm's \ac{EPLB}, a load-balancing  mechanism that periodically replicates and re-places experts across GPUs based on their observed load, so that each GPU receives a more even share of expert work.
We enable it on \textsc{Qwen3-30B-A3B} with 16 redundant experts.
Across our workloads, \ac{EPLB} reduces the fraction of prefill spent in bubbles from 11--17\% to 8--13\%.
Because \ac{EPLB} also changes how much decode and idle-time headroom serving leaves, we break down the contribution of placing fine-tuning in decode and idle time with \sys-\textsc{Flat}, and in the bubbles with \sys. 

\Cref{fig:eval-eplb} shows the results for 8k-token prompts at medium and high load.
At high load and across prompt lengths, bubble filling increases fine-tuning throughput over \sys-\textsc{Flat} by 12--23\% with \ac{EPLB}, compared to 27--37\% without it.
%Bubble filling therefore remains useful with \ac{EPLB}.
This is because \ac{EPLB} balances expert load averaged over recent steps, while the routing of each individual prefill still deviates from that average, so some GPUs still finish their experts early in every layer.
At high load, where decode and idle-time headroom is scarce, these residual bubbles are the main remaining source of fine-tuning capacity.
p90 \ac{TTFT} and \ac{TPOT} stay within 7\% of serving alone.

At medium load, bubble filling changes fine-tuning throughput by $-2$ to $+0.5$\%.
As expected, decode and idle-time headroom is plentiful at this load, and \sys-\textsc{Flat} already uses it.
The only workload where latency exceeds 7\% is also at medium load: at 16k-token prompts and 1 request/s, p90 \ac{TPOT} reaches 1.32$\times$ serving alone (Appendix~\ref{sec:appx-eplb}).
This suggests \sys should leave residual bubbles unused when conventional headroom already suffices.

Overall, \ac{EPLB} and bubble filling are complementary: \ac{EPLB} shrinks bubbles but cannot remove the imbalance within each prefill, and \sys turns what remains into 12--23\% more fine-tuning throughput at high load while protecting serving latency.
Bubble filling therefore pays off most when serving is busiest, which is exactly when conventional decode and idle-time headroom runs out.
